\section{Roadmap sections intentionally postponed}

This project is intended to become a reusable roadmap for similar tabular classification datasets. Some checks and methods are useful in general, but are intentionally postponed here because they are not necessary yet or depend on later modelling choices.

\subsection{Outlier diagnostics beyond visual inspection}

The numeric plots did not suggest an immediate need for outlier treatment. A future roadmap version may include an explicit IQR-based or robust-statistics outlier table:
\[
\text{IQR} = Q_3 - Q_1,
\]
\[
\text{lower bound} = Q_1 - 1.5 \cdot \text{IQR},
\qquad
\text{upper bound} = Q_3 + 1.5 \cdot \text{IQR}.
\]
This should be added only when it supports a decision. Outlier removal or capping should not be automatic.

\subsection{Missing-value mechanism analysis}

This dataset contains no remaining missing values after the confirmed \texttt{TotalCharges} correction. In future datasets, the report should include a more extensive missing-data mechanism analysis, distinguishing missing feature values from missing target labels and considering whether missing values are expected in production.

\subsection{Categorical association measures}

For categorical predictors, one could add Cramer's V, chi-square tests, or mutual information. These are not necessary for the current stage. The current report uses categorical frequency plots and churn-rate plots because they are easier to interpret and directly useful for exploratory understanding.

\subsection{Feature engineering}

Possible future feature engineering ideas include:
\begin{itemize}[leftmargin=*]
    \item grouping related service variables,
    \item counting the number of subscribed services,
    \item constructing tenure-charge ratios,
    \item representing contract length ordinally where justified.
\end{itemize}
These should be considered carefully because feature engineering choices may introduce leakage or impose assumptions.

\subsection{Model-specific preprocessing}

Model-specific preprocessing is postponed. Later scripts should build pipelines for each model family:
\begin{itemize}[leftmargin=*]
    \item linear models, kNN, SVMs, and neural networks typically require encoding and scaling;
    \item tree-based models require categorical encoding in scikit-learn but do not require scaling;
    \item Naive Bayes variants require special representations depending on whether Gaussian, Bernoulli, or multinomial assumptions are used.
\end{itemize}

\subsection{Validation strategy}

A fixed validation set was not created at this stage. Later model scripts can use cross-validation inside the training set, or create an internal validation split from the training set only. The held-out test set should remain unused until final evaluation.